% Compile with: xelatex -interaction=nonstopmode -halt-on-error report.tex
\documentclass[11pt,a4paper]{article}

\usepackage[vietnamese]{babel}
\usepackage{fontspec}
\usepackage[a4paper,margin=2cm]{geometry}
\usepackage{booktabs}
\usepackage{tabularx}
\usepackage{array}
\usepackage{enumitem}
\usepackage{xcolor}
\usepackage{titlesec}
\usepackage{microtype}
\usepackage{hyperref}

\setmainfont{DejaVu Serif}
\setsansfont{DejaVu Sans}
\setmonofont{Noto Sans Mono}
\newfontfamily\hanfont{Droid Sans Fallback}
\newcommand{\han}[1]{{\hanfont #1}}

% Điền thông tin nhóm trước khi nộp.
\newcommand{\reportauthor}{Nhóm sinh viên: \textit{[Điền họ tên và MSSV]}}
\newcommand{\coursename}{Học phần: \textit{[Điền tên học phần]}}

\definecolor{linkblue}{HTML}{174A7E}
\definecolor{bookcolor}{HTML}{DCEEFF}
\definecolor{timecolor}{HTML}{FFF0C2}
\definecolor{locationcolor}{HTML}{DDF4DF}
\definecolor{titlecolor}{HTML}{F3DDF5}
\hypersetup{
  colorlinks=true,
  linkcolor=linkblue,
  urlcolor=linkblue,
  pdftitle={Báo cáo xây dựng ngữ liệu đơn ngữ chữ Hán cổ},
  pdfauthor={Nhóm sinh viên}
}

\setlength{\parindent}{1.2em}
\setlength{\parskip}{0.25em}
\setlist{nosep,leftmargin=1.6em}
\titleformat{\section}{\large\bfseries}{\thesection.}{0.5em}{}
\titleformat{\subsection}{\normalsize\bfseries}{\thesubsection.}{0.5em}{}
\titlespacing*{\section}{0pt}{1.1em}{0.45em}
\titlespacing*{\subsection}{0pt}{0.8em}{0.3em}
\renewcommand{\arraystretch}{1.15}

\begin{document}

\begin{center}
  {\Large\bfseries BÁO CÁO XÂY DỰNG NGỮ LIỆU ĐƠN NGỮ\\
  CHỮ HÁN CỔ VỀ LỊCH SỬ TRUNG QUỐC}\par
  \vspace{0.6em}
  {\coursename}\par
  {\reportauthor}\par
  \vspace{0.35em}
  {\small Tháng 7 năm 2026}
\end{center}

\noindent\textbf{Tóm tắt.}
Báo cáo trình bày quy trình xây dựng ngữ liệu đơn ngữ chữ Hán cổ/văn ngôn thuộc
lĩnh vực lịch sử Trung Quốc thời phong kiến. Từ tám tác phẩm đầu vào, hệ thống
chuẩn hóa cấu trúc, làm sạch, tách câu và sinh tiền gán nhãn thực thể có tên
(NER). Kết quả hiện tại gồm 57\,474 bản ghi, 271\,471 câu và 192\,109 thực thể
tiền gán nhãn, được xuất thành 787 thư mục quyển. Đây là bản \emph{draft}: các
offset và cấu trúc đã vượt qua kiểm tra tự động, nhưng nhãn thực thể chưa được
xem là gold cho đến khi hoàn tất hiệu đính thủ công.

\section{Giới thiệu}

\subsection{Mục tiêu của đồ án}

Mục tiêu là tạo một corpus có cấu trúc phục vụ nghiên cứu xử lý ngôn ngữ tự
nhiên trên văn bản Hán cổ, tập trung vào sử liệu Trung Quốc thời phong kiến.
``Đơn ngữ'' được hiểu là trường văn bản nghiên cứu chỉ chứa tiếng Hán, không
ghép song song với bản dịch tiếng Việt. Thông tin tiếng Việt, nếu có, chỉ đóng
vai trò metadata nhận diện tác phẩm và không được đưa vào nội dung câu.

\subsection{Bài toán được giao}

Đầu vào là các tệp văn bản thô, trong đó tiêu đề tác phẩm, tiêu đề quyển, dòng
trống và nội dung chính còn trộn lẫn. Đồ án cần biến nguồn này thành dữ liệu
JSONL theo từng quyển, tách câu và gán các nhãn thực thể theo schema nộp bài:
\texttt{PER}, \texttt{LOC}, \texttt{ORG}, \texttt{TITLE}, \texttt{TME} và
\texttt{NUM}. Trong nội bộ, hệ thống giữ nhãn chi tiết hơn như
\texttt{BOOK}, \texttt{OFFICIAL\_TITLE}, \texttt{DYNASTY} và
\texttt{POLITY}, sau đó ánh xạ sang schema đầu ra.

\subsection{Phạm vi dữ liệu và kết quả mong đợi}

Phạm vi gồm tám tác phẩm có trong thư mục \path{dataset/}, không tự động tải
thêm nguồn bên ngoài và không thực hiện OCR. Sản phẩm mong đợi gồm corpus đã
làm sạch, corpus tiền gán nhãn, tập pilot phục vụ hiệu đính, bản xuất theo từng
quyển, manifest SHA-256 và báo cáo kiểm tra tuân thủ. Bản nộp cuối chỉ được tạo
khi mapping nhãn đã được xác nhận và toàn bộ dữ liệu trong phạm vi nộp đã được
con người duyệt.

\section{Dữ liệu và công cụ sử dụng}

\subsection{Mô tả và nguồn dữ liệu đầu vào}

Nguồn thô gồm tám tệp UTF-8 do đề bài cung cấp. Mỗi tệp tương ứng một tác phẩm;
các dòng dạng \han{卷一}, \han{卷上}, \han{卷}015b được dùng để xác định ranh
giới quyển. Danh sách nguồn được thể hiện trong Bảng~\ref{tab:sources}.

\begin{table}[ht]
\centering
\caption{Các tác phẩm trong dữ liệu đầu vào.}
\label{tab:sources}
\small
\begin{tabularx}{\textwidth}{>{\raggedright\arraybackslash}X l l}
\toprule
Tác phẩm & Giai đoạn & Thể loại/đặc trưng \\
\midrule
\han{北史} & Nam--Bắc triều & Chính sử \\
\han{北齊書} & Bắc Tề & Chính sử \\
\han{諸蕃志} & Tống & Địa lý, ngoại quốc \\
\han{舊唐書} & Đường & Chính sử \\
\han{舊五代史} & Ngũ Đại & Chính sử \\
\han{東觀漢記（四庫全書本）} & Đông Hán & Sử thư \\
\han{漢書} & Tây Hán & Chính sử \\
\han{後漢書} & Đông Hán & Chính sử \\
\bottomrule
\end{tabularx}
\end{table}

Quy tắc làm sạch bảo toàn nguyên văn: không dịch, không chuyển giản thể/phồn
thể và không tự sửa chữ nghi là lỗi OCR. Hệ thống chỉ loại tiêu đề công việc,
heading quyển, dòng trống và boilerplate đã nhận diện; mọi quyết định loại bỏ
được lưu trong \path{build/audit_decisions.jsonl}. Những đoạn chứa URL hoặc
placeholder OCR được giữ lại nhưng đánh dấu \texttt{needs\_review}.

\subsection{Công cụ, thư viện và mô hình}

Pipeline được viết bằng Python 3.11 trở lên, quản lý môi trường bằng \texttt{uv}
và cấu hình bằng TOML. Các thành phần chính gồm bộ tách câu dựa trên quy tắc,
biểu thức chính quy, từ điển dạng trie, seed gazetteer và bộ hợp nhất ứng viên
có thứ tự ưu tiên xác định. Dữ liệu duyệt được trao đổi theo định dạng tương
thích Doccano 1.8.x. Kiểm thử dùng \texttt{pytest}; kiểm tra mã nguồn dùng
\texttt{Ruff}.

Các adapter cho CBDB và CHGIS đã được chuẩn bị để tạo seed cache có manifest,
nhưng không tự tải dữ liệu và không tham gia nếu người chạy chưa cung cấp nguồn.
Các dependency cho mô hình phân đoạn/NER cũng được khai báo ở dạng tùy chọn.
Do đó, số liệu trong báo cáo này chỉ đến từ hai nguồn ứng viên
\texttt{regex} và \texttt{seed}, không phải dự đoán của mô hình học máy.

\section{Quy trình thực hiện}

\subsection{Tổng quan}

\begin{center}
\small
\fbox{Văn bản thô} $\rightarrow$
\fbox{Làm sạch và tách quyển} $\rightarrow$
\fbox{Tách câu} $\rightarrow$
\fbox{Sinh và hợp nhất ứng viên} $\rightarrow$
\fbox{Duyệt Doccano} $\rightarrow$
\fbox{Xuất và kiểm tra}
\end{center}

Quy trình tách rõ dữ liệu nguồn, tiền gán nhãn và gold. Cách tổ chức này tránh
việc xem kết quả heuristic là nhãn đúng, đồng thời cho phép truy vết từ thực thể
đầu ra về tệp và dòng nguồn.

\subsection{Các bước chính}

\begin{enumerate}
  \item \textbf{Chuẩn hóa và tách quyển.} Parser nhận diện nhiều biến thể heading,
  chuẩn hóa khóa \texttt{(volume\_number, volume\_part)} và kiểm tra đường dẫn
  đầu ra không va chạm.

  \item \textbf{Làm sạch có kiểm toán.} Những dòng bị loại được ghi cùng lý do;
  49 bản ghi có dấu hiệu URL/placeholder OCR được giữ lại để con người quyết định.

  \item \textbf{Tách câu.} Văn bản được chia theo \han{。！？}, đồng thời bảo vệ
  dấu câu nằm trong ngoặc. Mỗi câu có \texttt{sentence\_id}, offset bắt đầu/kết
  thúc và trạng thái duyệt.

  \item \textbf{Tiền gán nhãn.} Regex và seed gazetteer sinh ứng viên. Bộ merger
  xử lý cùng span, lồng nhau, giao cắt và tie-break một cách xác định. Mỗi thực
  thể giữ \texttt{sources}, \texttt{priority\_score} và danh sách nhãn đã hợp nhất.

  \item \textbf{Lấy mẫu và hiệu đính.} Tập pilot ngẫu nhiên gồm 200 câu, phân tầng
  đều theo tám tác phẩm (25 câu/tác phẩm, seed 42); 15\% được dành cho gán nhãn
  kép. Một tập diagnostic riêng tập trung vào các trường hợp khó.

  \item \textbf{Xuất dữ liệu.} Chế độ \texttt{draft} cho phép khảo sát tiền gán
  nhãn. Chế độ \texttt{final} bắt buộc strict validation, mapping được xác nhận,
  câu/thực thể đã duyệt và conflict đã giải quyết. Việc ghi output dùng staging
  và đổi tên trong cùng filesystem để tránh submission dở dang.
\end{enumerate}

\section{Kết quả đạt được}

\subsection{Khối lượng dữ liệu xử lý}

\begin{table}[ht]
\centering
\caption{Thống kê của lần chạy hiện tại.}
\label{tab:stats}
\small
\begin{tabular}{lr@{\qquad}lr}
\toprule
Chỉ tiêu & Số lượng & Chỉ tiêu & Số lượng \\
\midrule
Tệp/tác phẩm nguồn & 8 & Ký tự nội dung & 7\,359\,348 \\
Bản ghi & 57\,474 & Câu & 271\,471 \\
Quyển đầu ra & 787 & Thực thể tiền gán nhãn & 192\,109 \\
Dòng audit & 3\,817 & Bản ghi cần xem lại làm sạch & 49 \\
Conflict chưa giải quyết & 1\,362 & Bản ghi chứa conflict & 1\,224 \\
\bottomrule
\end{tabular}
\end{table}

Phân bố nhãn nội bộ được trình bày ở Bảng~\ref{tab:labels}. Hai nhãn
\texttt{PERSON} và \texttt{NUMBER} chưa có ứng viên trong lần chạy hiện tại;
đây là thiếu hụt coverage, không phải bằng chứng rằng corpus không có người hay
số lượng.

\begin{table}[ht]
\centering
\caption{Phân bố 192\,109 thực thể tiền gán nhãn.}
\label{tab:labels}
\small
\begin{tabular}{lr@{\qquad}lr}
\toprule
Nhãn & Số lượng & Nhãn & Số lượng \\
\midrule
\texttt{LOCATION} & 83\,061 & \texttt{OFFICIAL\_TITLE} & 63\,463 \\
\texttt{TIME} & 26\,035 & \texttt{BOOK} & 14\,267 \\
\texttt{DYNASTY} & 3\,034 & \texttt{POLITY} & 2\,249 \\
\bottomrule
\end{tabular}
\end{table}

\subsection{Các sản phẩm đầu ra}

\begin{table}[ht]
\centering
\small
\begin{tabularx}{\textwidth}{>{\ttfamily\raggedright\arraybackslash}p{0.42\textwidth} X}
\toprule
Artifact & Nội dung \\
\midrule
build/corpus.jsonl & Corpus đã chuẩn hóa và tách quyển/câu \\
build/corpus\_preannotated.jsonl & Corpus có thực thể heuristic và provenance \\
build/audit\_decisions.jsonl & Nhật ký quyết định làm sạch \\
build/pilot/ & Mẫu random và diagnostic để hiệu đính \\
build/submission/ & 787 thư mục quyển, manifest và validation report \\
build/annotation\_samples.html & Trang HTML minh họa trực quan các nhãn \\
\bottomrule
\end{tabularx}
\end{table}

Manifest của bản draft ghi 271\,471 câu, 192\,109 thực thể, mapping phiên bản
0.2.0 và checksum SHA-256 của corpus nguồn. Mỗi thực thể trong file nộp chỉ còn
bốn trường tối thiểu: \texttt{text}, \texttt{label}, \texttt{start},
\texttt{end}; provenance chỉ tồn tại trong artifact nội bộ.

\subsection{Ví dụ minh họa}

Ví dụ thực từ \han{北史}:

\begin{quote}
\han{\colorbox{timecolor}{二年}\,春正月癸卯，拜沮渠無諱為征西\colorbox{titlecolor}{大將軍}、
\colorbox{locationcolor}{涼州}\,牧、酒泉王。}
\end{quote}

Trong đó \han{二年} được gán \texttt{TIME}, \han{大將軍} là
\texttt{OFFICIAL\_TITLE}, và \han{涼州} là \texttt{LOCATION}. Ví dụ cũng cho
thấy hạn chế quan trọng: tên người \han{沮渠無諱} chưa được nhận diện và địa danh
\han{酒泉} có thể bị bỏ sót. Một ví dụ khác từ \han{諸蕃志} nhận diện
\han{寶慶元年} là \texttt{TIME} và \han{大夫} là
\texttt{OFFICIAL\_TITLE}. Vì đây là preannotation, các đoạn tô màu là đề xuất
cho người gán nhãn, không phải đáp án gold.

\section{Đánh giá và thảo luận}

\subsection{Cách đánh giá}

Đánh giá được chia thành hai tầng. Tầng cấu trúc kiểm tra offset, chuỗi
\texttt{text[start:end]}, nhãn hợp lệ, phạm vi câu, đường dẫn quyển, shape đầu
ra và chuỗi checksum. Tầng chất lượng NER dự kiến dùng tập gold pilot để tính
precision, recall và F1 strict theo entity, F1 biên, kết quả theo từng nhãn,
confusion matrix và độ chính xác nhãn trên các biên đã khớp.

\subsection{Kết quả đánh giá hiện tại}

\begin{table}[ht]
\centering
\small
\begin{tabularx}{\textwidth}{X l X}
\toprule
Kiểm tra & Kết quả & Diễn giải \\
\midrule
Offset-only validation & Đạt & Offset và cấu trúc cơ bản hợp lệ \\
Draft compliance & 0 fatal, 0 warning & Shape, thư mục và manifest hợp lệ \\
Kiểm thử phần mềm & 190 test đạt & Bảo vệ hành vi parser, merger, I/O và export \\
Ruff lint/format & Đạt & Mã nguồn vượt qua quality gate tĩnh \\
Strict validation preannotation & 464\,853 fatal & Kỳ vọng vì chưa human-review \\
Entity-level F1 & Chưa có & Chưa có gold pilot đã hiệu đính \\
\bottomrule
\end{tabularx}
\end{table}

Số lỗi strict gồm 271\,471 câu chưa duyệt, 192\,109 thực thể chưa duyệt,
1\,224 bản ghi còn conflict và 49 bản ghi cần quyết định làm sạch. Các số này
không nên được diễn giải là 464\,853 dự đoán sai; chúng là các cổng quy trình
ngăn bản preannotation bị xuất nhầm thành final. Tương tự, 190 kiểm thử phần
mềm chứng minh tính ổn định của implementation, không thay thế cho precision,
recall hoặc F1 của dữ liệu.

\subsection{Vấn đề gặp phải và phân tích nguyên nhân}

\begin{itemize}
  \item \textbf{Biến thể heading quyển.} Cùng một nguồn có nhiều cách ghi số và
  phần quyển. Parser ban đầu có thể gộp nhầm; nguyên nhân là xem chuỗi heading
  như tên tự do thay vì khóa chuẩn hóa.

  \item \textbf{Nhập nhằng thực thể.} Một span có thể là triều đại, chính thể,
  địa danh hoặc nằm trong chức danh dài. Regex chỉ quan sát bề mặt nên sinh cả
  false positive, false negative và conflict.

  \item \textbf{Coverage không cân bằng.} Gazetteer địa danh/chức quan phong phú
  hơn tên người và số lượng; đồng thời tên người Hán cổ thường không có dấu hiệu
  hình thức cố định. Đây là lý do \texttt{LOCATION} và
  \texttt{OFFICIAL\_TITLE} chiếm ưu thế, còn \texttt{PERSON}/\texttt{NUMBER}
  chưa được sinh trong lần chạy này.

  \item \textbf{Nhiễu nguồn.} Placeholder OCR và URL nằm giữa dòng không thể loại
  bằng quy tắc boilerplate mà không có nguy cơ xóa văn bản thật. Hệ thống vì vậy
  giữ nguyên và chuyển sang hàng đợi duyệt.

  \item \textbf{Thiếu gold.} Không thể báo cáo F1 đáng tin cậy bằng cách so
  heuristic với chính nó. Cần gán nhãn độc lập, giải quyết bất đồng và khóa phiên
  bản guideline trước khi đánh giá.
\end{itemize}

\section{Kết luận và hướng phát triển}

Đồ án đã hoàn thành đường ống có thể tái lập từ tám tệp nguồn đến corpus JSONL,
bao gồm tách 787 quyển, làm sạch có audit, tách 271\,471 câu, sinh 192\,109 thực
thể tiền gán nhãn, lấy mẫu pilot, trao đổi Doccano, export nguyên tử và kiểm tra
checksum. Source code vượt qua test và lint; bản draft vượt qua compliance.

Hạn chế lớn nhất là chưa có gold corpus đã được người làm chuyên môn duyệt,
mapping cuối chưa được xác nhận và coverage của \texttt{PERSON}/\texttt{NUMBER}
còn thiếu. Vì vậy, hướng phát triển ưu tiên là: (1) hiệu đính 200 câu pilot và
đo agreement giữa annotator; (2) phân tích lỗi theo nhãn để mở rộng gazetteer;
(3) tích hợp CBDB/CHGIS có kiểm soát; (4) thử mô hình pretrained cho Hán cổ và
so sánh với baseline regex--seed; (5) active learning để ưu tiên câu có conflict
hoặc độ bất định cao. Chỉ sau các bước này mới xác nhận mapping, chạy strict
validation và đóng gói bản final.

\begin{thebibliography}{9}
\small
\bibitem{evahan2024}
L. Li và cộng sự, ``EvaHan 2024: Overview of the First International Ancient
Chinese Translation Bakeoff,'' ACL Anthology, 2024.
\url{https://aclanthology.org/2024.lt4hala-1.15/}

\bibitem{evahan2025}
L. Li và cộng sự, ``EvaHan 2025: Overview of Named Entity Recognition and
Sentence Segmentation Evaluation for Ancient Chinese,'' ACL Anthology, 2025.
\url{https://aclanthology.org/2025.alp-1.17/}

\bibitem{guner2023}
Y. Xu và cộng sự, ``GuNER 2023: Ancient Chinese Named Entity Recognition
Evaluation,'' ACL Anthology, 2023.
\url{https://aclanthology.org/2023.lt4hala-1.18/}

\bibitem{datasetreadme}
Tài liệu nguồn của đồ án, \path{dataset/README.md} và các hướng dẫn trong
\path{docs/}, phiên bản truy cập tháng 7 năm 2026.
\end{thebibliography}

\end{document}
